\section{利用导数证明不等式}

\subsection{证明恒成立和存在问题}

\begin{theorem}[利用导数证明不等式]
    构造新函数，利用单调性和最值来证明不等式恒成立。
\end{theorem}
\begin{example}
    已知函数 $f(x) = \mathrm{e}^{x - 1} - 2x - 1$，其导函数为 $f'(x)$.
    \begin{enumerate}
        \item[(1)] 求 $f(x)$ 的单调区间；
        \item[(2)] 证明：$xf'(x) + x - \ln x \geq 0$.
    \end{enumerate}
\end{example}

\v{12em}

\subsection{含餐问题的解法}

\begin{theorem}[含参恒成立问题]
    \begin{itemize}
        \item \textbf{带着参数求最值}: 通常转化为求函数的最值问题。求出 $f(x)$ 的最小值，使其大于等于0，从而解出参数 $a$ 的范围。
        \item \textbf{参变分离求值域}: 通过固定参数 $a$，研究 $f(x)$ 的值域，进而确定 $a$ 的取值范围。
    \end{itemize}
\end{theorem}

\subsubsection{参变分离法}

\begin{example}
    已知函数 $f(x) = ax\ln x - 2x + a, a \in \mathbb{R}$.
    \begin{enumerate}
        \item[(1)] 若 $a = 1$，求曲线 $y = f(x)$ 在点 $(1, -1)$ 处的切线方程；
        \item[(2)] 若 $a > 0$，$f(x) \geq 0$ 恒成立，求 $a$ 的取值范围.
    \end{enumerate}
\end{example}

\subsubsection{带参运算法}
